% Recuperación literal de nucleo.tex REV08, líneas 10--32.
% Ampliación específica de II, situada después del núcleo común inalterado.
\subsection{Correspondencia entre soporte posicional y calendario orientado}
\label{ii:ampliacion-posicional-calendario}

La representación del calendario distingue las posiciones de APP y los
tránsitos orientados de su evolución. Retomamos esa correspondencia para
explicitar la lectura del segmento que contiene el empalme entre vueltas.

Consideremos el alfabeto de nueve índices
\[
 D_9=\{1,2,3,4,5,6,7,8,9\}.
\]
Su disposición horizontal y vertical produce las ochenta y una posiciones
de \(D_9^2\). Para describir el movimiento se utiliza además el ciclo
orientado \(C_9\). Sus nueve aristas reciben, en orden, los índices de
\(D_9\). El índice de una arista especifica un tránsito entre estados,
mientras que un vértice especifica su frontera. La aplicación que asigna
a cada arista su índice permite usar el mismo alfabeto en el soporte
posicional y en el calendario; conserva la distinción entre posición,
intervalo orientado y estado transportado.

En el levantamiento entero del calendario, la primera vuelta comprende
los nueve intervalos consecutivos \([0,1],\ldots,[8,9]\). La continuación
\([9,10]\) inicia la vuelta siguiente con el índice uno y con una memoria
actualizada. El segmento dibujado hasta diez muestra ambas vueltas en su
punto de empalme. Esta representación especifica los nueve tránsitos de
una vuelta y el siguiente tránsito; la holonomía del estado se construye
después mediante la composición efectiva del TPK.

\begin{figure}[htbp]
\centering
\begin{tikzpicture}[x=1.05cm,y=1cm,>=Latex,font=\small]
 \draw[very thick,ink] (0,0)--(9,0);
 \draw[very thick,dashed,accent] (9,0)--(10,0);
 \foreach \x in {0,...,10}{
   \draw (\x,-.08)--(\x,.08);
   \node[below=5pt,font=\scriptsize] at (\x,0) {$\x$};
 }
 \foreach \x in {1,...,9}{
  \draw[->,ink] (\x-.85,.35)--(\x-.15,.35);
  \node[above=3pt] at (\x-.5,.35) {$\x$};
 }
 \draw[->,accent] (9.15,.35)--(9.85,.35);
 \node[above=3pt,text=accent] at (9.5,.35) {$1$};
 \draw[decorate,decoration={brace,amplitude=5pt},ink] (0,1)--(9,1)
  node[midway,above=7pt] {Nueve intervalos orientados: una vuelta};
 \node[align=center] at (5,-1.05)
  {Índice de tránsito: $g(t)=1+(t\bmod9)$\\
   Memoria del calendario levantado: $m(t)=\lfloor t/9\rfloor$};
 \node[align=center] at (5,-1.85)
  {$g(t+9)=g(t),\qquad m(t+9)=m(t)+1$};
\end{tikzpicture}
\caption{Levantamiento del ciclo de nueve tránsitos. Los números superiores
etiquetan intervalos orientados; los inferiores localizan sus fronteras.
El intervalo discontinuo muestra el primer tránsito de la vuelta siguiente.
El índice de fase retorna y el cociente del calendario aumenta. Este esquema
representa el calendario; la acción sobre residuos, cocientes, orientación,
frontera y memoria se define mediante las operaciones del TPK.}
\label{fig:ii-segmento-generador}
\end{figure}

